\clearpage

\subsection{Dataset Acquisition}
Real-world datasets for Indore are compiled from multiple authoritative and
open-data sources to ensure spatial accuracy and modelling validity. The
dataset collection spans the following domains:

\begin{itemize}
    \item \textbf{Urban and demographic data:} ward-level population
          distribution, administrative ward boundaries, land-use
          classification, and residential/commercial density across
          the 85~wards of Indore Municipal Corporation.
    \item \textbf{Candidate venue data:} locations and attributes of
          parking lots, CNG/fuel stations, restaurants, and shopping
          malls identified as permissible hosts for public charging
          infrastructure under MoRTH~(2022) guidelines, along with
          venue-type-specific charger count capacity and estimated
          capital expenditure.
    \item \textbf{EV adoption and charging behaviour:} current EV
          penetration statistics, projected logistic adoption
          trajectories, typical charging-session duration (30~minutes),
          and public charging fraction ($F_{\text{public}} = 0.20$).
    \item \textbf{Power-grid infrastructure:} zone-level substation
          capacities and grid upgrade cost estimates used to enforce
          per-zone grid capacity constraints in the optimization loop.
    \item \textbf{Economic and tariff parameters:} venue-type
          installation costs (₹6.5--18.0~lakh), Time-of-Use electricity
          tariffs (₹112--₹210 per session across Idle, Normal, and
          Peak periods), and operating cost factors for annual
          profitability computation.
\end{itemize}

\noindent\textbf{Purpose:} These datasets collectively define ward-level
demand patterns, the heterogeneous candidate venue pool, grid constraints,
and the economic context for the four-objective optimization.

%---------------------------------------------------------
\subsection{Data Preprocessing and Spatial Zoning}
The raw datasets are cleaned, normalised, and transformed into coherent
spatial and numerical formats suitable for optimization and queueing
evaluation.

\begin{itemize}
    \item Partition Indore's 85~wards into five functional zones ---
          Central Core, North, South, East, and West --- based on
          geographic location and land-use characteristics, assigning
          each zone a population estimate, projected EV count, daily
          charging demand, and zone-adjusted service radius
          ($R_z = 1.5$~km for Central Core; $R_z = 2.5$~km for all
          Outer zones).
    \item Collate the base pool of $30+$ candidate stations across
          four venue types and augment it dynamically: wards above
          the 75th~percentile of projected demand receive up to
          3~synthetic parking candidates; wards above the 50th
          percentile receive up to 2, simulating municipal land
          acquisition at an elevated base cost of ₹7.0~lakh.
    \item Compute a ward--station Euclidean distance matrix to
          support the spatial coverage objective and zone-radius
          constraint checks.
    \item Compute the per-station base daily demand $\mu_j$ by
          allocating ward-level demand $D_w(t)$ proportionally among
          candidate stations within the ward's service radius.
    \item Filter and validate all entries, removing inconsistencies,
          duplicates, and physically infeasible values (negative
          distances, zero capacities, out-of-bound costs).
\end{itemize}

\noindent\textbf{Purpose:} Preprocessing ensures spatial integrity and
produces all derived inputs --- distance matrix, demand estimates,
zone assignments, and cost vectors --- required by the optimization
and queueing stages.

%---------------------------------------------------------
\subsection{Demand and Scenario Modelling}
EV charging demand is projected across the planning horizon using a
two-stage demographic and adoption model, supplemented by stochastic
simulation:

\begin{itemize}
    \item Ward-level population is forecast to the target year using
          a constant CAGR of $r = 1.8\%$ applied to the 2011~Census
          baseline:
          $\text{Pop}_w(t) = \text{Pop}_w(2011)\times(1.018)^{t-2011}$.
    \item The aggregate EV penetration rate follows a logistic
          S-curve with parameters $K = 1.0$, $a = 0.3$, and
          inflection year $t_0 = 2030$, consistent with NITI~Aayog
          targets and battery cost parity trajectories.
    \item Daily public charging demand per ward is estimated as:
          $D_w(t) = \text{Pop}_w(t)\times(V/1000)\times E(t)
          \times F_{\text{public}}$, where $V = 50$~EVs per
          1000~persons and $F_{\text{public}} = 0.20$.
    \item Three planning scenarios (2025 baseline, 2030 mid-term,
          2035 long-term) are generated by evaluating the logistic
          curve at the respective years, with the 2025~scenario
          serving as the primary optimization target.
    \item A 100-scenario Monte~Carlo stochastic overlay is applied
          per candidate station: $D_{j,s} \sim \max(0,\,
          \mathcal{N}(\mu_j,\,(v\mu_j)^2))$, $s = 1,\dots,100$,
          with noise fraction $v = 0.15$, ensuring robustness
          against demand edge cases.
\end{itemize}

\noindent\textbf{Purpose:} This two-stage pipeline produces
deterministic ward-level demand forecasts validated by the logistic
curve parameterisation, and stochastic scenario envelopes that
prevent over-optimistic station configurations from being selected
by the evolutionary optimizers.

%---------------------------------------------------------
\subsection{Optimization Model Formulation}
The EVCS placement problem is formulated as a four-objective binary
optimisation problem (MINLP).

\textbf{Decision variable:}
\begin{itemize}
    \item Binary station-selection vector
          $\mathbf{x} = [x_1,\dots,x_N]^{\top}$,
          $x_j \in \{0,1\}$, indicating whether candidate
          station $j$ is deployed.
\end{itemize}

\textbf{Objectives} (all cast as minimisation):
\begin{itemize}
    \item $f_1$: Minimize total capital expenditure (installation
          and grid upgrade costs).
    \item $f_2$: Maximize spatial demand coverage (negated fraction
          of total daily EV sessions covered within zone-adjusted
          service radii).
    \item $f_3$: Maximize net return on investment (negated ratio
          of annual profit to CAPEX, replacing raw profit to avoid
          Pareto front degeneracy).
    \item $f_4$: Minimize mean queue waiting time across all selected
          stations, evaluated via the M/M/$c$/K model under
          Monte~Carlo demand averaging.
\end{itemize}

\textbf{Constraints:}
\begin{itemize}
    \item Per-zone grid capacity limits on aggregated charger load,
    \item Zone-adjusted service radius coverage requirement
          ($R_z = 1.5$~km / $2.5$~km),
    \item Hard blocking probability constraint: $P_{\text{block}}
          \leq 10\%$ per station, enforced as a $+50$-minute
          penalty on $f_4$,
    \item Queue stability: utilisation $\rho_j < 1$ at every
          selected station.
\end{itemize}

\noindent\textbf{Purpose:} This formulation integrates siting,
service quality, and financial viability into a unified Pareto
optimisation framework without requiring pre-specified objective
weights.

%---------------------------------------------------------
\subsection{Multi-Objective Optimization: NSGA-II and MOPSO}
Both algorithms are applied independently to the formulated problem
and their Pareto fronts are compared using formal performance
indicators.

\begin{itemize}
    \item \textbf{NSGA-II:} A population of 80~binary individuals
          is evolved over up to 500~generations using uniform
          crossover ($p_c = 0.9$) and bit-flip mutation
          ($p_m = 1/n$, which flips one bit per genome on average).
          Fast non-dominated sorting and
          crowding distance selection maintain a well-distributed
          Pareto front. Elitism prevents loss of good solutions
          across generations.
    \item \textbf{MOPSO:} A swarm of 80~particles navigates a
          continuous $[0,1]^N$ position space, binarised at
          evaluation time via a 0.5~threshold. An external archive
          of up to 200~non-dominated solutions is maintained;
          the global guide is selected probabilistically, biased
          toward less-crowded archive regions. Constriction
          weights $w = 0.4$ and coefficients $c_1 = c_2 = 1.5$
          govern swarming behaviour.
    \item \textbf{Early stopping:} Both algorithms halt if the
          Hypervolume improvement falls below $\Delta\text{HV}
          < 0.0005$ over 40~consecutive generations/iterations,
          with a minimum guaranteed run of 100~generations.
    \item \textbf{Performance indicators:} Hypervolume~(HV),
          Generational Distance~(GD), and Spread are computed
          over the four-dimensional objective space to formally
          quantify solution quality and diversity.
\end{itemize}

\noindent\textbf{Purpose:} Running both algorithms independently
with identical evaluation budgets and comparing their outputs via
formal metrics yields an objective, reproducible assessment of
which algorithm better suits the binary EVCS placement problem
structure.

%---------------------------------------------------------
\subsection{Queueing-Based Operational Evaluation}
Queueing theory is applied within the evolutionary evaluation loop
to verify operational feasibility and enforce service quality
constraints.

\begin{itemize}
    \item The daily demand $\mu_j$ (averaged over 100~Monte~Carlo
          scenarios) is converted to an hourly arrival rate
          $\lambda_j$ and further scaled by the ToU demand
          multiplier for each period (Peak $\times 2.8$,
          Normal $\times 1.0$, Idle $\times 0.15$).
    \item Each selected station is modelled as an M/M/$c_j$/K
          finite-capacity queue with $K = c_j + 8$ physical
          waiting bays, mean service rate $\mu = 2$~sessions/hour
          (30-minute sessions), and FCFS discipline.
    \item Steady-state performance metrics are computed analytically
          for each ToU period:
          \begin{itemize}
              \item blocking probability $P_{\text{block},j}$,
              \item mean queue waiting time $W_{q,j}$,
              \item server utilisation $\rho_j = \lambda_j / (c_j\mu)$,
              \item expected queue length $L_{q,j}$.
          \end{itemize}
    \item The reported $\overline{W}_{q,j}$ is the time-weighted
          average across all three ToU periods over the 24-hour
          cycle.
    \item Stations with $P_{\text{block}} > 10\%$ incur a
          $+50$-minute penalty on objective $f_4$, and stations
          with $\rho_j \geq 1$ are flagged as unstable and
          penalised accordingly.
\end{itemize}

\noindent\textbf{Purpose:} Embedding the M/M/$c$/K evaluation
inside the optimization loop --- rather than applying it
post-hoc --- forces both algorithms to actively avoid
queue-infeasible configurations, producing deployment plans
that are operationally stable under peak demand conditions.

%---------------------------------------------------------
\subsection{Sensitivity Analysis and Scalability Benchmarking}
The robustness and computational viability of the framework are
validated through two additional evaluation stages.

\begin{itemize}
    \item \textbf{Sensitivity analysis:} The impact of four
          parameters on objective values is examined independently:
          \begin{itemize}
              \item EV growth rate (logistic curve steepness $a$
                    scaled $\pm 20\%$),
              \item electricity cost (commercial grid tariff
                    varied by $\pm$₹2.0/kWh),
              \item installation cost factor (base CAPEX scaled
                    $\pm 15\%$),
              \item charger service rate (mean session duration
                    varied from 20 to 45~minutes).
          \end{itemize}
          For each perturbation, both algorithms are re-run and the
          shift in the optimal solution's $f_1$--$f_4$ values is
          recorded.
    \item \textbf{Scalability benchmarking:} The candidate pool
          size is varied from $N = 5$ to $N = 20$ stations.
          Median wall-clock runtime is measured over 5~independent
          trials for both NSGA-II and MOPSO, and a polynomial
          scaling curve is fitted to confirm tractability for
          larger metropolitan deployments.
\end{itemize}

\noindent\textbf{Purpose:} Sensitivity analysis quantifies the
resilience of the recommended deployment plan to macroeconomic
and technological shifts, while scalability benchmarks confirm
that the Python-based framework is applicable beyond the Indore
case study without requiring high-performance computing resources.

%---------------------------------------------------------
\subsection{Results Aggregation, Visualization, and Interpretation}
Outputs from the optimization, queueing, and sensitivity stages
are synthesised into visual and tabular summaries.

\begin{itemize}
    \item Generate individual Pareto front plots for NSGA-II and
          MOPSO across all six pairs of the four objectives, and
          a combined comparison plot identifying regions of
          algorithmic dominance.
    \item Plot HV convergence curves per generation/iteration for
          both algorithms, annotating the fastest-gain generation
          and the HV plateau point.
    \item Visualize the queue waiting time distribution across
          all Pareto solutions for both algorithms, highlighting
          the 20-minute user-tolerance threshold.
    \item Summarize algorithm performance in a structured table
          reporting HV, GD, Spread, median $W_q$, max coverage,
          and max net ROI for both NSGA-II and MOPSO.
    \item Report the recommended optimal solution (maximum
          coverage subject to $P_{\text{block}} \leq 10\%$) with
          full metrics: stations selected, total CAPEX, coverage,
          net ROI, and mean $W_q$.
    \item Formulate phased deployment recommendations and identify
          wards requiring priority station placement or grid
          reinforcement under the 2030 and 2035 demand scenarios.
\end{itemize}

\noindent\textbf{Purpose:} This step transforms raw Pareto front
data and queueing metrics into deployment-ready planning insights
accessible to municipal authorities, infrastructure investors,
and urban mobility planners.